\documentclass[a4paper,12pt,twocolumn]{article}

\usepackage[landscape,top=1cm,bottom=1.4cm,left=1cm,right=1cm,columnsep=1cm]{geometry}
\usepackage[fleqn]{amsmath}
\usepackage{amssymb}
\setlength{\mathindent}{1.5em}
\usepackage{fontspec}
\usepackage{enumitem}
\usepackage{framed}
\usepackage{needspace}
\usepackage{xeCJK}

% Set fonts
\setmainfont{Times New Roman}
\setsansfont{Arial}
\setmonofont{JetBrains Mono}

% Chinese font support
\setCJKmainfont{Iansui}[
  BoldFont=Iansui,
  AutoFakeBold=3.17
]

\pagestyle{plain}
\setlength{\columnseprule}{0.2pt}
\newcommand{\examdate}{2026/03/11}
\newenvironment{solutionbox}{\begin{framed}}{\end{framed}}

\begin{document}

\twocolumn[
{\centering
\Large\bfseries Calculus Problems Solutions\par
\vspace{0.5em}
\noindent\hfill\normalsize\normalfont \examdate\par
\vspace{0.8em}
}
]

\begin{enumerate}[label=\textbf{\arabic*.}, leftmargin=2em, itemsep=1.2em]
\setcounter{enumi}{48}

\Needspace{28\baselineskip}
\item Determine whether the series is convergent or divergent. If it is convergent, find its sum:
\[
\sum_{n=1}^{\infty}\left(\frac{1}{e^n}+\frac{1}{n(n+1)}\right)
\]

\begin{solutionbox}
\textbf{Method 1: Split the series into two known series}

Write
\[
\sum_{n=1}^{\infty}\left(\frac{1}{e^n}+\frac{1}{n(n+1)}\right)
=
\sum_{n=1}^{\infty}\frac{1}{e^n}
+
\sum_{n=1}^{\infty}\frac{1}{n(n+1)}.
\]
For the first series,
\[
\sum_{n=1}^{\infty}\frac{1}{e^n}
=
\sum_{n=1}^{\infty}\left(\frac{1}{e}\right)^n
=
\frac{1/e}{1-1/e}
=
\frac{1}{e-1}.
\]
For the second series, use partial fractions:
\[
\frac{1}{n(n+1)}=\frac{1}{n}-\frac{1}{n+1}.
\]
So
\[
\sum_{n=1}^{\infty}\frac{1}{n(n+1)}
=
\sum_{n=1}^{\infty}\left(\frac{1}{n}-\frac{1}{n+1}\right).
\]
This is a telescoping series. Its \(N\)-th partial sum is
\[
S_N=
\left(1-\frac12\right)+\left(\frac12-\frac13\right)+\cdots+\left(\frac1N-\frac1{N+1}\right)
=
1-\frac{1}{N+1}.
\]
Taking \(N\to\infty\), we get
\[
\sum_{n=1}^{\infty}\frac{1}{n(n+1)}=1.
\]
Therefore,
\[
\sum_{n=1}^{\infty}\left(\frac{1}{e^n}+\frac{1}{n(n+1)}\right)
=
\frac{1}{e-1}+1
=
\frac{e}{e-1}.
\]
\end{solutionbox}

\begin{solutionbox}
\textbf{Method 2: Work with partial sums directly}

Let
\[
S_N=\sum_{n=1}^{N}\left(\frac{1}{e^n}+\frac{1}{n(n+1)}\right).
\]
Then
\[
S_N=
\sum_{n=1}^{N}\frac{1}{e^n}
+
\sum_{n=1}^{N}\frac{1}{n(n+1)}.
\]
For the finite geometric sum,
\[
\sum_{n=1}^{N}\frac{1}{e^n}
=
\sum_{n=1}^{N}\left(\frac{1}{e}\right)^n
=
\frac{\frac1e\left(1-\left(\frac1e\right)^N\right)}{1-\frac1e}
=
\frac{1-(1/e)^N}{e-1}.
\]
For the second part,
\[
\frac{1}{n(n+1)}=\frac{1}{n}-\frac{1}{n+1},
\]
so
\[
\sum_{n=1}^{N}\frac{1}{n(n+1)}
=
\sum_{n=1}^{N}\left(\frac{1}{n}-\frac{1}{n+1}\right)
=
1-\frac{1}{N+1}.
\]
Hence
\[
S_N=
\frac{1-(1/e)^N}{e-1}+1-\frac{1}{N+1}.
\]
Now let \(N\to\infty\). Since
\[
\left(\frac{1}{e}\right)^N\to 0
\qquad\text{and}\qquad
\frac{1}{N+1}\to 0,
\]
we obtain
\[
\lim_{N\to\infty}S_N
=
\frac{1}{e-1}+1
=
\frac{e}{e-1}.
\]

\medskip
\textbf{Answer:}
\[
\boxed{\sum_{n=1}^{\infty}\left(\frac{1}{e^n}+\frac{1}{n(n+1)}\right)\text{ converges, and its sum is }\frac{e}{e-1}.}
\]
\end{solutionbox}

\Needspace{20\baselineskip}
\item Determine whether the series is convergent or divergent. If it is convergent, find its sum:
\[
\sum_{n=1}^{\infty}\frac{e^n}{n^2}
\]

\begin{solutionbox}
\textbf{Method 1: nth-term test with L'Hôpital's Rule}

Let
\[
a_n=\frac{e^n}{n^2}.
\]
A necessary condition for the convergence of \(\sum a_n\) is
\[
\lim_{n\to\infty} a_n = 0.
\]
So we examine the limit of \(a_n\).

Consider the reciprocal:
\[
\lim_{n\to\infty}\frac{n^2}{e^n}
=
\lim_{x\to\infty}\frac{x^2}{e^x}.
\]
This is an \(\infty/\infty\) form, so apply L'Hôpital's Rule twice:
\[
\lim_{x\to\infty}\frac{x^2}{e^x}
=
\lim_{x\to\infty}\frac{2x}{e^x}
=
\lim_{x\to\infty}\frac{2}{e^x}
=0.
\]
Therefore,
\[
\lim_{n\to\infty}\frac{n^2}{e^n}=0
\quad\Longrightarrow\quad
\lim_{n\to\infty}\frac{e^n}{n^2}=\infty.
\]
Hence \(a_n \not\to 0\). Since the terms of the series do not approach \(0\), the series diverges.
\end{solutionbox}

\begin{solutionbox}
\textbf{Method 2: Exponential growth dominates polynomial growth}

Let
\[
a_n=\frac{e^n}{n^2}.
\]
The numerator \(e^n\) has exponential growth, while the denominator \(n^2\) has polynomial growth. Exponential growth is much faster than polynomial growth, so
\[
\frac{e^n}{n^2}\to\infty
\qquad (n\to\infty).
\]
In particular,
\[
\lim_{n\to\infty}\frac{e^n}{n^2}\neq 0.
\]
Therefore the necessary condition for convergence fails, so the series diverges.

\end{solutionbox}

\setcounter{enumi}{76}
\Needspace{24\baselineskip}
\item Find the value of \(c\) if
\[
\sum_{n=2}^{\infty}(1+c)^{-n}=2
\]

\begin{solutionbox}
\textbf{Method 1: Treat it as a geometric series directly}

Let
\[
x=1+c.
\]
Then
\[
\sum_{n=2}^{\infty}(1+c)^{-n}
=
\sum_{n=2}^{\infty}x^{-n}
\]
is a geometric series with first term \(a=x^{-2}\) and common ratio \(r=x^{-1}\). Hence
\[
\frac{x^{-2}}{1-x^{-1}}=2.
\]
Simplify:
\[
\frac{1/x^2}{1-1/x}
=
\frac{1}{x(x-1)}
=2.
\]
So
\[
1=2x(x-1)
\quad\Longrightarrow\quad
2x^2-2x-1=0.
\]
Using the quadratic formula,
\[
x=\frac{2\pm\sqrt{(-2)^2-4(2)(-1)}}{4}
=
\frac{1\pm\sqrt{3}}{2}.
\]
For the infinite geometric series to converge, we need
\[
|r|<1
\quad\Longleftrightarrow\quad
\left|\frac{1}{x}\right|<1
\quad\Longleftrightarrow\quad
|x|>1.
\]
Thus the valid value is
\[
x=\frac{1+\sqrt{3}}{2}.
\]
Therefore,
\[
c=x-1=\frac{1+\sqrt{3}}{2}-1=\frac{\sqrt{3}-1}{2}.
\]
\end{solutionbox}

\begin{solutionbox}
\textbf{Method 2: Substitute \(r=(1+c)^{-1}\)}

Let
\[
r=(1+c)^{-1}.
\]
Then
\[
\sum_{n=2}^{\infty}(1+c)^{-n}
=
\sum_{n=2}^{\infty}r^n
=
\frac{r^2}{1-r}=2,
\]
with convergence condition \(|r|<1\).

So
\[
r^2=2(1-r)
\quad\Longrightarrow\quad
r^2+2r-2=0.
\]
Hence
\[
r=\frac{-2\pm\sqrt{4+8}}{2}
=
-1\pm\sqrt{3}.
\]
Since \(|r|<1\), the valid value is
\[
r=\sqrt{3}-1.
\]
Now
\[
1+c=\frac{1}{r}=\frac{1}{\sqrt{3}-1}
=
\frac{\sqrt{3}+1}{2},
\]
so
\[
c=\frac{\sqrt{3}+1}{2}-1=\frac{\sqrt{3}-1}{2}.
\]

\medskip
\textbf{Answer:}
\[
\boxed{c=\frac{\sqrt{3}-1}{2}}
\]
\end{solutionbox}

\setcounter{enumi}{7}
\Needspace{28\baselineskip}
\item Use the Integral Test to determine whether the series is convergent or divergent:
\[
\sum_{n=1}^{\infty} n^2 e^{-n^3}
\]

\begin{solutionbox}
\textbf{Method 1: Integral Test}

Let
\[
f(x)=x^2e^{-x^3}, \qquad x\ge 1.
\]
For \(x\ge 1\), \(f(x)>0\) and \(f(x)\) is continuous. To use the Integral Test, we also check that \(f(x)\) is decreasing.

Differentiate:
\[
f'(x)=\frac{d}{dx}\left(x^2e^{-x^3}\right)
=2xe^{-x^3}+x^2e^{-x^3}(-3x^2).
\]
So
\[
f'(x)=e^{-x^3}(2x-3x^4)=xe^{-x^3}(2-3x^3).
\]
For \(x\ge 1\), we have \(2-3x^3<0\), while \(x>0\) and \(e^{-x^3}>0\). Hence
\[
f'(x)<0 \qquad \text{for } x\ge 1,
\]
so \(f(x)\) is decreasing.

Now evaluate the improper integral:
\[
\int_1^\infty x^2e^{-x^3}\,dx.
\]
Use the substitution \(u=x^3\), so \(du=3x^2\,dx\) and \(x^2\,dx=\frac13\,du\). Then
\[
\int_1^\infty x^2e^{-x^3}\,dx
=
\frac13\int_1^\infty e^{-u}\,du
=
\frac13\left[-e^{-u}\right]_1^\infty
=
\frac{1}{3e}.
\]
Since this improper integral is finite, the Integral Test shows that
\[
\sum_{n=1}^{\infty} n^2 e^{-n^3}
\]
converges.
\end{solutionbox}

\begin{solutionbox}
\textbf{Method 2: Comparison with a convergent geometric series}

For sufficiently large \(n\), we have
\[
n^2\le e^n.
\]
Therefore,
\[
n^2e^{-n^3}\le e^n e^{-n^3}=e^{\,n-n^3}.
\]
Since \(n^3-n\ge n\) for all sufficiently large \(n\),
\[
e^{\,n-n^3}\le e^{-n}.
\]
Hence, for all sufficiently large \(n\),
\[
0\le n^2e^{-n^3}\le e^{-n}.
\]
But
\[
\sum_{n=1}^{\infty} e^{-n}
\]
is a geometric series with ratio \(1/e<1\), so it converges. By the Comparison Test,
\[
\sum_{n=1}^{\infty} n^2 e^{-n^3}
\]
also converges.

\medskip
\textbf{Answer:}
\[
\boxed{\sum_{n=1}^{\infty} n^2 e^{-n^3}\text{ converges}.}
\]
\end{solutionbox}

\setcounter{enumi}{9}
\Needspace{30\baselineskip}
\item Use the Integral Test to determine whether the series is convergent or divergent:
\[
\sum_{n=1}^{\infty}\frac{\tan^{-1} n}{1+n^2}
\]

\begin{solutionbox}
\textbf{Method 1: Integral Test}

Let
\[
f(x)=\frac{\tan^{-1}x}{1+x^2}, \qquad x\ge 1.
\]
For \(x\ge 1\), \(f(x)>0\) and \(f(x)\) is continuous. To use the Integral Test, we also check that \(f(x)\) is decreasing.

Differentiate:
\[
f'(x)=\frac{(1+x^2)\cdot \frac{1}{1+x^2}-(\tan^{-1}x)(2x)}{(1+x^2)^2}
=\frac{1-2x\tan^{-1}x}{(1+x^2)^2}.
\]
For \(x\ge 1\),
\[
\tan^{-1}x\ge \tan^{-1}1=\frac{\pi}{4},
\]
so
\[
2x\tan^{-1}x\ge 2\cdot 1\cdot \frac{\pi}{4}=\frac{\pi}{2}>1.
\]
Hence
\[
1-2x\tan^{-1}x<0.
\]
Since the denominator is always positive,
\[
f'(x)<0 \qquad \text{for } x\ge 1,
\]
so \(f(x)\) is decreasing.

Now evaluate the improper integral:
\[
\int_1^\infty \frac{\tan^{-1}x}{1+x^2}\,dx.
\]
Use the substitution
\[
u=\tan^{-1}x, \qquad du=\frac{1}{1+x^2}\,dx.
\]
Then
\[
\int_1^\infty \frac{\tan^{-1}x}{1+x^2}\,dx
=
\int_{\pi/4}^{\pi/2} u\,du
=
\left[\frac{u^2}{2}\right]_{\pi/4}^{\pi/2}
=
\frac{3\pi^2}{32}.
\]
Since this improper integral is finite, the Integral Test shows that
\[
\sum_{n=1}^{\infty}\frac{\tan^{-1} n}{1+n^2}
\]
converges.
\end{solutionbox}

\begin{solutionbox}
\textbf{Method 2: Comparison Test}

For every \(n\ge 1\),
\[
\tan^{-1}n<\frac{\pi}{2}.
\]
Therefore,
\[
0\le \frac{\tan^{-1}n}{1+n^2}
\le \frac{\pi/2}{1+n^2}.
\]
Also, since \(1+n^2\ge n^2\),
\[
\frac{1}{1+n^2}\le \frac{1}{n^2}.
\]
Hence
\[
0\le \frac{\tan^{-1}n}{1+n^2}
\le \frac{\pi}{2n^2}.
\]
Now
\[
\sum_{n=1}^{\infty}\frac{\pi}{2n^2}
=
\frac{\pi}{2}\sum_{n=1}^{\infty}\frac{1}{n^2}
\]
converges because \(\sum 1/n^2\) is a \(p\)-series with \(p=2>1\). By the Comparison Test,
\[
\sum_{n=1}^{\infty}\frac{\tan^{-1} n}{1+n^2}
\]
also converges.

\medskip
\textbf{Answer:}
\[
\boxed{\sum_{n=1}^{\infty}\frac{\tan^{-1} n}{1+n^2}\text{ converges}.}
\]
\end{solutionbox}

\setcounter{enumi}{31}
\Needspace{34\baselineskip}
\item Find the values of \(p\) for which the series is convergent:
\[
\sum_{n=3}^{\infty} \frac{1}{n \ln n \,[\ln(\ln n)]^p}
\]

\begin{solutionbox}
\textbf{Method 1: Integral Test}

Let
\[
f(x)=\frac{1}{x\ln x\,[\ln(\ln x)]^p}, \qquad x\ge 3.
\]
For \(x\ge 3\), we have \(\ln x>0\) and \(\ln(\ln x)>0\), so \(f(x)>0\), and \(f(x)\) is continuous. It is also eventually decreasing, so the Integral Test applies.

Consider
\[
\int_3^\infty \frac{dx}{x\ln x\,[\ln(\ln x)]^p}.
\]
Use the substitution
\[
u=\ln(\ln x), \qquad du=\frac{dx}{x\ln x}.
\]
Then the integral becomes
\[
\int u^{-p}\,du.
\]

If \(p\neq 1\), then
\[
\int u^{-p}\,du=\frac{u^{1-p}}{1-p},
\]
which converges at \(\infty\) exactly when \(1-p<0\), that is, when \(p>1\).

If \(p=1\), then
\[
\int u^{-1}\,du=\ln u,
\]
which diverges.

Therefore the improper integral converges exactly when \(p>1\). By the Integral Test, the series converges exactly when
\[
p>1.
\]
\end{solutionbox}

\begin{solutionbox}
\textbf{Method 2: Cauchy Condensation Test}

Let
\[
a_n=\frac{1}{n\ln n\,[\ln(\ln n)]^p}.
\]
Since \(a_n>0\) and is eventually decreasing, Cauchy Condensation gives
\[
\sum a_n \text{ converges }
\Longleftrightarrow
\sum 2^k a_{2^k} \text{ converges.}
\]
Now
\[
2^k a_{2^k}
=
2^k\cdot \frac{1}{2^k\ln(2^k)\,[\ln(\ln(2^k))]^p}
=
\frac{1}{k\ln 2\,[\ln(k\ln 2)]^p}.
\]
This has the same convergence behavior as
\[
\sum_{k=2}^{\infty}\frac{1}{k(\ln k)^p},
\]
because
\[
\lim_{k\to\infty}\left(\frac{\ln k}{\ln(k\ln 2)}\right)^p=1.
\]

Apply Cauchy Condensation again to
\[
b_k=\frac{1}{k(\ln k)^p}.
\]
Then
\[
2^m b_{2^m}
=
2^m\cdot \frac{1}{2^m(\ln(2^m))^p}
=
\frac{1}{(m\ln 2)^p}.
\]
So \(\sum b_k\) has the same convergence behavior as
\[
\sum_{m=1}^{\infty}\frac{1}{m^p},
\]
which is a \(p\)-series and converges exactly when \(p>1\).

\medskip
\textbf{Answer:}
\[
\boxed{\text{The series converges exactly for } p>1.}
\]
\end{solutionbox}

\setcounter{enumi}{33}
\Needspace{34\baselineskip}
\item Find the values of \(p\) for which the series is convergent:
\[
\sum_{n=1}^{\infty}\frac{\ln n}{n^p}
\]

\begin{solutionbox}
\textbf{Method 1: Integral Test}

Let
\[
f(x)=\frac{\ln x}{x^p}, \qquad x\ge 1.
\]
If \(p\le 0\), then
\[
a_n=\frac{\ln n}{n^p}
\]
does not tend to \(0\), so the series diverges immediately.

Now assume \(p>0\). Then \(f(x)\) is positive and continuous for \(x\ge 1\). Also,
\[
f'(x)=x^{-p-1}(1-p\ln x),
\]
so \(f'(x)<0\) for sufficiently large \(x\). Hence \(f(x)\) is eventually decreasing, and the Integral Test applies.

Consider
\[
\int_1^\infty \frac{\ln x}{x^p}\,dx.
\]

If \(p\neq 1\), use integration by parts with
\[
u=\ln x, \qquad dv=x^{-p}\,dx.
\]
Then
\[
\int \frac{\ln x}{x^p}\,dx
=
\frac{x^{1-p}\ln x}{1-p}-\frac{x^{1-p}}{(1-p)^2}+C.
\]
If \(p>1\), then \(x^{1-p}\to 0\) and \(x^{1-p}\ln x\to 0\), so the improper integral converges.
If \(0<p<1\), then \(x^{1-p}\to\infty\), so the improper integral diverges.

If \(p=1\), then
\[
\int_1^\infty \frac{\ln x}{x}\,dx
=
\left[\frac{(\ln x)^2}{2}\right]_1^\infty
=\infty,
\]
so it diverges.

Therefore, by the Integral Test, the series converges exactly when
\[
p>1.
\]
\end{solutionbox}

\begin{solutionbox}
\textbf{Method 2: Comparison with a \(p\)-series}

First consider \(p\le 1\). For \(n\ge 3\), we have \(\ln n\ge 1\), so
\[
\frac{\ln n}{n^p}\ge \frac{1}{n^p}.
\]
But
\[
\sum_{n=3}^{\infty}\frac{1}{n^p}
\]
diverges for \(p\le 1\). Hence
\[
\sum_{n=1}^{\infty}\frac{\ln n}{n^p}
\]
also diverges for \(p\le 1\).

Now assume \(p>1\). Choose
\[
\varepsilon=\frac{p-1}{2}>0.
\]
Since \(\ln n<n^\varepsilon\) for all sufficiently large \(n\), we get
\[
\frac{\ln n}{n^p}
\le
\frac{n^\varepsilon}{n^p}
=
\frac{1}{n^{p-\varepsilon}}
=
\frac{1}{n^{(p+1)/2}}
\]
for all sufficiently large \(n\).

Because
\[
\frac{p+1}{2}>1,
\]
the series
\[
\sum_{n=1}^{\infty}\frac{1}{n^{(p+1)/2}}
\]
converges. By the Comparison Test,
\[
\sum_{n=1}^{\infty}\frac{\ln n}{n^p}
\]
also converges for \(p>1\).

\medskip
\textbf{Answer:}
\[
\boxed{\text{The series converges exactly when } p>1.}
\]
\end{solutionbox}

\setcounter{enumi}{45}
\Needspace{34\baselineskip}
\item For what values of \(p\) does the series
\[
\sum_{n=2}^{\infty}\frac{1}{n^p\ln n}
\]
converge?

\begin{solutionbox}
\textbf{Method 1: Integral Test}

Let
\[
f(x)=\frac{1}{x^p\ln x}, \qquad x\ge 2.
\]
If \(p<0\), then
\[
\frac{1}{n^p\ln n}=\frac{n^{-p}}{\ln n},
\]
which does not tend to \(0\), so the series diverges immediately.

Now assume \(p\ge 0\). Then \(f(x)\) is positive and decreasing on \([2,\infty)\), so the Integral Test applies. Consider
\[
\int_2^\infty \frac{dx}{x^p\ln x}.
\]
Use the substitution
\[
t=\ln x, \qquad x=e^t,\qquad dx=e^t\,dt,\qquad x^p=e^{pt}.
\]
Then
\[
\int_2^\infty \frac{dx}{x^p\ln x}
=
\int_{\ln 2}^{\infty}\frac{e^{(1-p)t}}{t}\,dt.
\]

If \(p>1\), then \(1-p<0\), so \(e^{(1-p)t}\) decays exponentially. Hence the integral converges.

If \(p=1\), then the integral becomes
\[
\int_{\ln 2}^{\infty}\frac{dt}{t},
\]
which diverges.

If \(0\le p<1\), then \(1-p>0\), so \(e^{(1-p)t}\) grows exponentially, and for sufficiently large \(t\),
\[
\frac{e^{(1-p)t}}{t}\ge \frac{1}{t}.
\]
Since \(\int_{\ln 2}^{\infty}\frac{dt}{t}\) diverges, the integral also diverges.

Therefore, by the Integral Test, the series converges exactly when
\[
p>1.
\]
\end{solutionbox}

\begin{solutionbox}
\textbf{Method 2: Cauchy Condensation Test}

For \(p\ge 0\), let
\[
a_n=\frac{1}{n^p\ln n}.
\]
Then \(a_n\) is positive and decreasing, so Cauchy Condensation applies:
\[
\sum_{n=2}^{\infty} a_n \text{ converges }
\Longleftrightarrow
\sum_{k=1}^{\infty}2^k a_{2^k}\text{ converges.}
\]
Now
\[
2^k a_{2^k}
=
2^k\cdot \frac{1}{(2^k)^p\ln(2^k)}
=
\frac{2^{k(1-p)}}{k\ln 2}.
\]
So we only need to study
\[
\sum_{k=1}^{\infty}\frac{2^{k(1-p)}}{k}.
\]

If \(p>1\), then \(1-p<0\), so \(2^{k(1-p)}\) decays geometrically and
\[
\frac{2^{k(1-p)}}{k}\le 2^{k(1-p)}.
\]
Hence the condensed series converges.

If \(p=1\), then the condensed series becomes a constant multiple of
\[
\sum_{k=1}^{\infty}\frac{1}{k},
\]
so it diverges.

If \(0\le p<1\), then \(1-p>0\), so
\[
\frac{2^{k(1-p)}}{k}
\]
does not tend to \(0\), and the condensed series diverges.

If \(p<0\), the original terms do not tend to \(0\), so the original series diverges immediately.

\medskip
\textbf{Answer:}
\[
\boxed{\text{The series converges exactly when } p>1.}
\]
\end{solutionbox}

\setcounter{enumi}{40}
\Needspace{32\baselineskip}
\item Approximate the sum of the series correct to four decimal places:
\[
\sum_{n=1}^{\infty}\frac{(-1)^n}{(2n)!}
\]

\begin{solutionbox}
\textbf{Method 1: Recognize the cosine series}

Recall the Maclaurin series
\[
\cos x=\sum_{n=0}^{\infty}\frac{(-1)^n x^{2n}}{(2n)!}.
\]
Setting \(x=1\), we get
\[
\cos 1=\sum_{n=0}^{\infty}\frac{(-1)^n}{(2n)!}.
\]
Our series starts at \(n=1\), so we subtract the \(n=0\) term:
\[
\sum_{n=1}^{\infty}\frac{(-1)^n}{(2n)!}
=
\cos 1-1.
\]
Since
\[
\cos 1\approx 0.5403023059,
\]
we obtain
\[
\cos 1-1\approx -0.4596976941.
\]
Therefore, correct to four decimal places,
\[
\boxed{-0.4597}.
\]
\end{solutionbox}

\begin{solutionbox}
\textbf{Method 2: Alternating Series Estimation Theorem}

Let
\[
a_n=\frac{1}{(2n)!}.
\]
Then the series is
\[
\sum_{n=1}^{\infty}(-1)^n a_n.
\]
The terms \(a_n\) are positive, decreasing, and tend to \(0\), so the Alternating Series Estimation Theorem applies.

Compute the first few partial sums:
\[
S_1=-\frac{1}{2!}=-0.5,
\]
\[
S_2=-\frac{1}{2!}+\frac{1}{4!}
=-0.5+\frac{1}{24}
=-0.4583333333,
\]
\[
S_3=-\frac{1}{2!}+\frac{1}{4!}-\frac{1}{6!}
=-0.4597222222.
\]
The next term has size
\[
\frac{1}{8!}=\frac{1}{40320}\approx 0.0000248016.
\]
Since
\[
0.0000248016<0.00005,
\]
the error in using \(S_3\) is less than \(0.00005\), which is enough for four-decimal accuracy. Thus
\[
\sum_{n=1}^{\infty}\frac{(-1)^n}{(2n)!}\approx -0.4597222222,
\]
so to four decimal places,
\[
\boxed{-0.4597}.
\]
\end{solutionbox}

\setcounter{enumi}{42}
\Needspace{32\baselineskip}
\item Approximate the sum of the series correct to four decimal places:
\[
\sum_{n=1}^{\infty} (-1)^n n e^{-2n}
\]

\begin{solutionbox}
\textbf{Method 1: Use the formula for \(\sum nr^n\)}

Let
\[
r=-e^{-2}.
\]
Then
\[
\sum_{n=1}^{\infty} (-1)^n n e^{-2n}
=
\sum_{n=1}^{\infty} nr^n.
\]
Since
\[
|r|=e^{-2}<1,
\]
we may use
\[
\sum_{n=1}^{\infty} nr^n=\frac{r}{(1-r)^2}.
\]
Therefore,
\[
\sum_{n=1}^{\infty} (-1)^n n e^{-2n}
=
\frac{-e^{-2}}{(1+e^{-2})^2}.
\]
Multiply numerator and denominator by \(e^4\):
\[
\frac{-e^{-2}}{(1+e^{-2})^2}
=
-\frac{e^2}{(e^2+1)^2}.
\]
Thus the exact sum is
\[
-\frac{e^2}{(e^2+1)^2}.
\]
Numerically,
\[
-\frac{e^2}{(e^2+1)^2}\approx -0.1049935854,
\]
so correct to four decimal places,
\[
\boxed{-0.1050}.
\]
\end{solutionbox}

\begin{solutionbox}
\textbf{Method 2: Derive the formula from a geometric series}

Let
\[
S=\sum_{n=1}^{\infty} nr^n,
\qquad \text{where } r=-e^{-2}.
\]
Since \(|r|<1\), the geometric series converges:
\[
r+r^2+r^3+\cdots=\frac{r}{1-r}.
\]
Now write
\[
S=r+2r^2+3r^3+4r^4+\cdots
\]
and multiply by \(r\):
\[
rS=r^2+2r^3+3r^4+4r^5+\cdots.
\]
Subtract:
\[
S-rS=r+r^2+r^3+r^4+\cdots=\frac{r}{1-r}.
\]
Hence
\[
S(1-r)=\frac{r}{1-r}
\quad\Longrightarrow\quad
S=\frac{r}{(1-r)^2}.
\]
Substitute \(r=-e^{-2}\):
\[
S=\frac{-e^{-2}}{(1+e^{-2})^2}
=
-\frac{e^2}{(e^2+1)^2}
\approx -0.1049935854.
\]

\medskip
\textbf{Answer:}
\[
\boxed{-0.1050}
\]
\end{solutionbox}

\setcounter{enumi}{39}
\Needspace{48\baselineskip}
\item The Bessel function of order 1 is defined by
\[
J_1(x)=\sum_{n=0}^{\infty}\frac{(-1)^n x^{2n+1}}{n!(n+1)!2^{2n+1}}
\]
\begin{enumerate}[label=(\alph*), leftmargin=2em]
\item Find the domain of \(J_1\).
\item Show that \(J_1\) satisfies the differential equation
\[
x^2 J_1''(x)+x J_1'(x)+(x^2-1)J_1(x)=0
\]
\item Show that \(J_0'(x)=-J_1(x)\).
\end{enumerate}

\begin{solutionbox}
\textbf{Method 1: Direct power-series calculation}

\textbf{(a)} Let
\[
u_n(x)=\frac{(-1)^n x^{2n+1}}{n!(n+1)!2^{2n+1}}.
\]
Apply the Ratio Test:
\[
\left|\frac{u_{n+1}(x)}{u_n(x)}\right|
=
\frac{|x|^2}{4(n+1)(n+2)}\to 0
\qquad (n\to\infty).
\]
So the defining series converges for every real \(x\). Therefore
\[
\boxed{\text{Domain of }J_1=(-\infty,\infty).}
\]

\textbf{(b)} Write
\[
J_1(x)=\sum_{n=0}^{\infty}a_nx^{2n+1},
\qquad
a_n=\frac{(-1)^n}{n!(n+1)!2^{2n+1}}.
\]
Since the radius of convergence is infinite, term-by-term differentiation is valid:
\[
J_1'(x)=\sum_{n=0}^{\infty}(2n+1)a_nx^{2n},
\qquad
J_1''(x)=\sum_{n=1}^{\infty}(2n+1)(2n)a_nx^{2n-1}.
\]
Hence
\[
x^2J_1''+xJ_1'-J_1
=
\sum_{n=0}^{\infty}4n(n+1)a_nx^{2n+1},
\]
and
\[
x^2J_1
=
\sum_{n=0}^{\infty}a_nx^{2n+3}
=
\sum_{n=1}^{\infty}a_{n-1}x^{2n+1}.
\]
So
\[
x^2J_1''+xJ_1'+(x^2-1)J_1
=
\sum_{n=1}^{\infty}\bigl(4n(n+1)a_n+a_{n-1}\bigr)x^{2n+1}.
\]
Now
\[
4n(n+1)a_n
=
\frac{(-1)^n}{(n-1)!n!2^{2n-1}},
\qquad
a_{n-1}
=
\frac{(-1)^{n-1}}{(n-1)!n!2^{2n-1}},
\]
so \(4n(n+1)a_n+a_{n-1}=0\) for all \(n\ge 1\). Therefore
\[
\boxed{x^2J_1''(x)+xJ_1'(x)+(x^2-1)J_1(x)=0.}
\]

\textbf{(c)} Use
\[
J_0(x)=\sum_{n=0}^{\infty}\frac{(-1)^n x^{2n}}{(n!)^2 2^{2n}}.
\]
Differentiating term by term,
\[
J_0'(x)=\sum_{n=1}^{\infty}\frac{(-1)^n(2n)x^{2n-1}}{(n!)^2 2^{2n}}
=
\sum_{n=1}^{\infty}\frac{(-1)^n x^{2n-1}}{(n-1)!n!2^{2n-1}}.
\]
Let \(m=n-1\). Then
\[
J_0'(x)
=
\sum_{m=0}^{\infty}\frac{(-1)^{m+1}x^{2m+1}}{m!(m+1)!2^{2m+1}}
=
-\sum_{m=0}^{\infty}\frac{(-1)^m x^{2m+1}}{m!(m+1)!2^{2m+1}}
=
-J_1(x).
\]
Hence
\[
\boxed{J_0'(x)=-J_1(x).}
\]
\end{solutionbox}

\begin{solutionbox}
\textbf{Method 2: Coefficient/recurrence viewpoint}

\textbf{(a)} For any real \(x\),
\[
\left|\frac{(-1)^n x^{2n+1}}{n!(n+1)!2^{2n+1}}\right|
=
|x|\cdot \frac{(|x|^2/4)^n}{n!(n+1)!}
\le
|x|\cdot \frac{(|x|^2/4)^n}{n!}.
\]
The comparison series
\[
\sum_{n=0}^{\infty}|x|\cdot \frac{(|x|^2/4)^n}{n!}
=
|x|e^{|x|^2/4}
\]
converges, so \(J_1\) converges absolutely for every real \(x\). Thus
\[
\boxed{\text{Domain of }J_1=(-\infty,\infty).}
\]

\textbf{(b)} Again write
\[
J_1(x)=\sum_{n=0}^{\infty}a_nx^{2n+1},
\qquad
a_n=\frac{(-1)^n}{n!(n+1)!2^{2n+1}}.
\]
Then
\[
x^2J_1''+xJ_1'+(x^2-1)J_1
=
\sum_{n=1}^{\infty}\bigl(4n(n+1)a_n+a_{n-1}\bigr)x^{2n+1}.
\]
So it is enough to verify the recurrence
\[
a_n=-\frac{a_{n-1}}{4n(n+1)}.
\]
Indeed,
\[
-\frac{a_{n-1}}{4n(n+1)}
=
-\frac{1}{4n(n+1)}
\cdot
\frac{(-1)^{n-1}}{(n-1)!n!2^{2n-1}}
=
\frac{(-1)^n}{n!(n+1)!2^{2n+1}}
=
a_n.
\]
Hence every coefficient vanishes, so
\[
\boxed{x^2J_1''(x)+xJ_1'(x)+(x^2-1)J_1(x)=0.}
\]

\textbf{(c)} Write
\[
J_0(x)=\sum_{n=0}^{\infty}b_nx^{2n},
\qquad
b_n=\frac{(-1)^n}{(n!)^2 2^{2n}}.
\]
Then
\[
J_0'(x)=\sum_{n=0}^{\infty}2(n+1)b_{n+1}x^{2n+1}.
\]
Now
\[
2(n+1)b_{n+1}
=
2(n+1)\cdot \frac{(-1)^{n+1}}{((n+1)!)^2 2^{2n+2}}
=
\frac{(-1)^{n+1}}{n!(n+1)!2^{2n+1}}
=
-a_n.
\]
Therefore
\[
J_0'(x)=\sum_{n=0}^{\infty}(-a_n)x^{2n+1}=-J_1(x).
\]

\medskip
\textbf{Answer:}
\[
\boxed{\text{Domain of }J_1=(-\infty,\infty),\quad
x^2J_1''+xJ_1'+(x^2-1)J_1=0,\quad
J_0'(x)=-J_1(x).}
\]
\end{solutionbox}

\setcounter{enumi}{48}
\Needspace{36\baselineskip}
\item Use the power series for \(\tan^{-1}x\) to prove the following expression for \(\pi\) as the sum of an infinite series:
\[
\pi = 2\sqrt{3}\sum_{n=0}^{\infty}\frac{(-1)^n}{(2n+1)3^n}
\]

\begin{solutionbox}
\textbf{Method 1: Direct substitution into the power series for \(\tan^{-1}x\)}

We use
\[
\tan^{-1}x=\sum_{n=0}^{\infty}\frac{(-1)^n x^{2n+1}}{2n+1},
\qquad |x|<1.
\]
Choose
\[
x=\frac{1}{\sqrt{3}},
\]
since
\[
\tan\left(\frac{\pi}{6}\right)=\frac{1}{\sqrt{3}}
\quad\Longrightarrow\quad
\tan^{-1}\left(\frac{1}{\sqrt{3}}\right)=\frac{\pi}{6}.
\]
Because \(\left|\frac{1}{\sqrt{3}}\right|<1\), substitution is valid:
\[
\tan^{-1}\left(\frac{1}{\sqrt{3}}\right)
=
\sum_{n=0}^{\infty}\frac{(-1)^n}{2n+1}\left(\frac{1}{\sqrt{3}}\right)^{2n+1}.
\]
Now
\[
\left(\frac{1}{\sqrt{3}}\right)^{2n+1}
=
\frac{1}{3^n\sqrt{3}},
\]
so
\[
\frac{\pi}{6}
=
\tan^{-1}\left(\frac{1}{\sqrt{3}}\right)
=
\frac{1}{\sqrt{3}}\sum_{n=0}^{\infty}\frac{(-1)^n}{(2n+1)3^n}.
\]
Multiplying both sides by \(6\), we get
\[
\pi
=
\frac{6}{\sqrt{3}}\sum_{n=0}^{\infty}\frac{(-1)^n}{(2n+1)3^n}
=
2\sqrt{3}\sum_{n=0}^{\infty}\frac{(-1)^n}{(2n+1)3^n}.
\]
\end{solutionbox}

\begin{solutionbox}
\textbf{Method 2: Derive the arctangent series from the geometric series}

Start from
\[
\frac{1}{1-r}=\sum_{n=0}^{\infty}r^n,
\qquad |r|<1.
\]
Let \(r=-x^2\). Then for \(|x|<1\),
\[
\frac{1}{1+x^2}
=
\sum_{n=0}^{\infty}(-1)^n x^{2n}.
\]
Integrate term by term from \(0\) to \(x\):
\[
\int_0^x \frac{1}{1+t^2}\,dt
=
\sum_{n=0}^{\infty}(-1)^n\int_0^x t^{2n}\,dt.
\]
Hence
\[
\tan^{-1}x
=
\sum_{n=0}^{\infty}\frac{(-1)^n x^{2n+1}}{2n+1}.
\]
Now substitute \(x=\frac{1}{\sqrt{3}}\):
\[
\tan^{-1}\left(\frac{1}{\sqrt{3}}\right)
=
\sum_{n=0}^{\infty}\frac{(-1)^n}{2n+1}\left(\frac{1}{\sqrt{3}}\right)^{2n+1}
=
\frac{1}{\sqrt{3}}\sum_{n=0}^{\infty}\frac{(-1)^n}{(2n+1)3^n}.
\]
Since
\[
\tan^{-1}\left(\frac{1}{\sqrt{3}}\right)=\frac{\pi}{6},
\]
we obtain
\[
\frac{\pi}{6}
=
\frac{1}{\sqrt{3}}\sum_{n=0}^{\infty}\frac{(-1)^n}{(2n+1)3^n}.
\]
Multiplying by \(6\),
\[
\pi
=
2\sqrt{3}\sum_{n=0}^{\infty}\frac{(-1)^n}{(2n+1)3^n}.
\]

\medskip
\textbf{Answer:}
\[
\boxed{\pi = 2\sqrt{3}\sum_{n=0}^{\infty}\frac{(-1)^n}{(2n+1)3^n}.}
\]
\end{solutionbox}

\setcounter{enumi}{37}
\Needspace{50\baselineskip}
\item The period of a pendulum with length \(L\) that makes a maximum angle \(\theta_0\) with the vertical is
\[
T = 4\sqrt{\frac{L}{g}} \int_0^{\pi/2} \frac{dx}{\sqrt{1-k^2\sin^2 x}},
\qquad
k=\sin\left(\frac{\theta_0}{2}\right).
\]
\begin{enumerate}[label=(\alph*), leftmargin=2em]
\item Expand the integrand as a binomial series and use the result of Exercise 7.1.56 to show that
\[
T = 2\pi\sqrt{\frac{L}{g}}
\left[
1+\frac{1^2}{2^2}k^2+\frac{1^2 3^2}{2^2 4^2}k^4+\frac{1^2 3^2 5^2}{2^2 4^2 6^2}k^6+\cdots
\right]
\]
\item Show that
\[
2\pi\sqrt{\frac{L}{g}}\left(1+\frac{1}{4}k^2\right)
\leq
T
\leq
2\pi\sqrt{\frac{L}{g}}\frac{4-3k^2}{4-4k^2}
\]
\item Use the inequalities in part (b) to estimate the period of a pendulum with \(L=1\) meter and \(\theta_0=10^\circ\). How does it compare with
\[
T \approx 2\pi\sqrt{\frac{L}{g}}\,?
\]
What if \(\theta_0=42^\circ\)?
\end{enumerate}

\begin{solutionbox}
\textbf{Method 1: Direct binomial-series expansion}

\textbf{(a)} Write
\[
\frac{1}{\sqrt{1-k^2\sin^2 x}}=(1-k^2\sin^2 x)^{-1/2}.
\]
Using the binomial series
\[
(1-u)^{-1/2}
=
1+\frac12u+\frac{1\cdot 3}{2\cdot 4}u^2+\frac{1\cdot 3\cdot 5}{2\cdot 4\cdot 6}u^3+\cdots,
\]
with \(u=k^2\sin^2 x\), we get
\[
\frac{1}{\sqrt{1-k^2\sin^2 x}}
=
1+\frac12k^2\sin^2x+\frac{1\cdot 3}{2\cdot 4}k^4\sin^4x+\frac{1\cdot 3\cdot 5}{2\cdot 4\cdot 6}k^6\sin^6x+\cdots.
\]
Integrating term by term from \(0\) to \(\pi/2\),
\[
\int_0^{\pi/2}\frac{dx}{\sqrt{1-k^2\sin^2x}}
=
\int_0^{\pi/2}dx
+\frac12k^2\int_0^{\pi/2}\sin^2x\,dx
+\frac{1\cdot 3}{2\cdot 4}k^4\int_0^{\pi/2}\sin^4x\,dx+\cdots.
\]
By Exercise 7.1.56,
\[
\int_0^{\pi/2}\sin^{2n}x\,dx
=
\frac{1\cdot 3\cdot 5\cdots(2n-1)}{2\cdot 4\cdot 6\cdots 2n}\cdot\frac{\pi}{2}.
\]
Therefore
\[
\int_0^{\pi/2}\frac{dx}{\sqrt{1-k^2\sin^2x}}
=
\frac{\pi}{2}
\left[
1+\frac{1^2}{2^2}k^2+\frac{1^2 3^2}{2^2 4^2}k^4+\frac{1^2 3^2 5^2}{2^2 4^2 6^2}k^6+\cdots
\right].
\]
Multiplying by \(4\sqrt{L/g}\),
\[
\boxed{
T=
2\pi\sqrt{\frac{L}{g}}
\left[
1+\frac{1^2}{2^2}k^2+\frac{1^2 3^2}{2^2 4^2}k^4+\frac{1^2 3^2 5^2}{2^2 4^2 6^2}k^6+\cdots
\right].
}
\]
Hence
\[
T\approx 2\pi\sqrt{\frac{L}{g}}
\]
using one term, and
\[
T\approx 2\pi\sqrt{\frac{L}{g}}\left(1+\frac14k^2\right)
\]
using two terms.

\textbf{(b)} All coefficients after the first are nonnegative, and for \(n\ge 1\),
\[
\left(
\frac{1\cdot 3\cdot 5\cdots (2n-1)}{2\cdot 4\cdot 6\cdots 2n}
\right)^2
\le \frac14.
\]
So
\[
T\ge 2\pi\sqrt{\frac{L}{g}}\left(1+\frac14k^2\right).
\]
Also,
\[
T\le
2\pi\sqrt{\frac{L}{g}}
\left[
1+\frac14k^2+\frac14k^4+\frac14k^6+\cdots
\right].
\]
Since
\[
k^2+k^4+k^6+\cdots=\frac{k^2}{1-k^2},
\]
we obtain
\[
T\le
2\pi\sqrt{\frac{L}{g}}
\left(1+\frac{k^2}{4(1-k^2)}\right)
=
2\pi\sqrt{\frac{L}{g}}\frac{4-3k^2}{4-4k^2}.
\]
Thus
\[
\boxed{
2\pi\sqrt{\frac{L}{g}}\left(1+\frac14k^2\right)
\le T\le
2\pi\sqrt{\frac{L}{g}}\frac{4-3k^2}{4-4k^2}.
}
\]

\textbf{(c)} Take \(L=1\) and \(g=9.8\). Then
\[
2\pi\sqrt{\frac{1}{9.8}}\approx 2.0061\text{ s}.
\]
If \(\theta_0=10^\circ\), then
\[
k=\sin 5^\circ\approx 0.08716,\qquad k^2\approx 0.007596.
\]
So
\[
2.0061\left(1+\frac14k^2\right)\approx 2.0099\text{ s},
\]
and
\[
2.0061\cdot \frac{4-3k^2}{4-4k^2}\approx 2.0099\text{ s}.
\]
Thus \(T\approx 2.0099\) s, which differs from the small-angle estimate \(2.0061\) s by only about \(0.0038\) s.

If \(\theta_0=42^\circ\), then
\[
k=\sin 21^\circ\approx 0.35837,\qquad k^2\approx 0.12843.
\]
So
\[
2.0061\left(1+\frac14k^2\right)\approx 2.0705\text{ s},
\]
and
\[
2.0061\cdot \frac{4-3k^2}{4-4k^2}\approx 2.0800\text{ s}.
\]
Hence
\[
\boxed{2.0705\text{ s}\lesssim T\lesssim 2.0800\text{ s}.}
\]
Compared with \(2.0061\) s, the small-angle approximation is now noticeably low.
\end{solutionbox}

\begin{solutionbox}
\textbf{Method 2: Coefficient-pattern method}

\textbf{(a)} Start from the generalized binomial theorem:
\[
(1-u)^{-1/2}
=
\sum_{n=0}^{\infty}\binom{-1/2}{n}(-u)^n.
\]
With \(u=k^2\sin^2x\),
\[
\frac{1}{\sqrt{1-k^2\sin^2x}}
=
\sum_{n=0}^{\infty}\binom{-1/2}{n}(-1)^n k^{2n}\sin^{2n}x.
\]
Now
\[
\binom{-1/2}{n}(-1)^n
=
\frac{1\cdot 3\cdot 5\cdots(2n-1)}{2\cdot 4\cdot 6\cdots 2n}.
\]
So
\[
\frac{1}{\sqrt{1-k^2\sin^2x}}
=
\sum_{n=0}^{\infty}
\frac{1\cdot 3\cdot 5\cdots(2n-1)}{2\cdot 4\cdot 6\cdots 2n}
k^{2n}\sin^{2n}x.
\]
Integrating term by term and using Exercise 7.1.56,
\[
\int_0^{\pi/2}\frac{dx}{\sqrt{1-k^2\sin^2x}}
=
\frac{\pi}{2}\sum_{n=0}^{\infty}
\left(
\frac{1\cdot 3\cdot 5\cdots(2n-1)}{2\cdot 4\cdot 6\cdots 2n}
\right)^2
k^{2n}.
\]
Thus
\[
\boxed{
T=
2\pi\sqrt{\frac{L}{g}}
\sum_{n=0}^{\infty}
\left(
\frac{1\cdot 3\cdot 5\cdots(2n-1)}{2\cdot 4\cdot 6\cdots 2n}
\right)^2
k^{2n}.
}
\]
Writing out the first terms gives the required series in part (a).

\textbf{(b)} Let
\[
c_n=
\left(
\frac{1\cdot 3\cdot 5\cdots(2n-1)}{2\cdot 4\cdot 6\cdots 2n}
\right)^2.
\]
Then \(c_1=\frac14\), and because each factor \(\frac{2j-1}{2j}<1\), we have
\[
c_n\le \frac14\qquad (n\ge 1).
\]
Therefore
\[
T=
2\pi\sqrt{\frac{L}{g}}
\left(1+c_1k^2+c_2k^4+\cdots\right)
\ge
2\pi\sqrt{\frac{L}{g}}\left(1+\frac14k^2\right),
\]
and also
\[
T\le
2\pi\sqrt{\frac{L}{g}}
\left(1+\frac14k^2+\frac14k^4+\frac14k^6+\cdots\right)
=
2\pi\sqrt{\frac{L}{g}}\frac{4-3k^2}{4-4k^2}.
\]
So the same inequality from part (b) follows.

\textbf{(c)} The numerical estimates are therefore the same as in Method 1:
\[
T\approx 2.0099\text{ s}\qquad (\theta_0=10^\circ),
\]
and
\[
2.0705\text{ s}\lesssim T\lesssim 2.0800\text{ s}
\qquad (\theta_0=42^\circ).
\]

\medskip
\textbf{Answer:}
\[
\boxed{
T=
2\pi\sqrt{\frac{L}{g}}
\left[
1+\frac{1^2}{2^2}k^2+\frac{1^2 3^2}{2^2 4^2}k^4+\cdots
\right]
}
\]
and
\[
\boxed{
2\pi\sqrt{\frac{L}{g}}\left(1+\frac14k^2\right)
\le T\le
2\pi\sqrt{\frac{L}{g}}\frac{4-3k^2}{4-4k^2}.
}
\]
For \(L=1\), this gives \(T\approx 2.0099\) s when \(\theta_0=10^\circ\), and
\[
2.0705\text{ s}\lesssim T\lesssim 2.0800\text{ s}
\]
when \(\theta_0=42^\circ\).
\end{solutionbox}

\setcounter{enumi}{2}
\Needspace{26\baselineskip}
\item Find an equation of sphere with center \((-3,2,5)\) and radius \(4\). What is the intersection of this sphere with the \(yz\)-plane?

\begin{solutionbox}
\textbf{Method 1: Use the standard equation of a sphere}

A sphere with center \((a,b,c)\) and radius \(r\) has equation
\[
(x-a)^2+(y-b)^2+(z-c)^2=r^2.
\]
Here the center is \((-3,2,5)\) and the radius is \(4\), so
\[
(x+3)^2+(y-2)^2+(z-5)^2=16.
\]
This is the equation of the sphere.

To find the intersection with the \(yz\)-plane, set
\[
x=0.
\]
Then
\[
(0+3)^2+(y-2)^2+(z-5)^2=16,
\]
so
\[
9+(y-2)^2+(z-5)^2=16
\]
and therefore
\[
(y-2)^2+(z-5)^2=7.
\]
Thus the intersection is a circle in the plane \(x=0\) with center \((0,2,5)\) and radius \(\sqrt{7}\).
\end{solutionbox}

\begin{solutionbox}
\textbf{Method 2: Use the geometry of a plane slice}

The sphere has center \((-3,2,5)\) and radius \(4\), so its equation is
\[
(x+3)^2+(y-2)^2+(z-5)^2=16.
\]
The \(yz\)-plane is the plane
\[
x=0.
\]
The distance from the center \((-3,2,5)\) to this plane is
\[
|-3|=3.
\]
When a plane cuts a sphere of radius \(r\) at distance \(d\) from its center, the cross-section is a circle of radius
\[
\sqrt{r^2-d^2}.
\]
So here the radius of the intersection circle is
\[
\sqrt{4^2-3^2}=\sqrt{16-9}=\sqrt{7}.
\]
Its center is the projection of \((-3,2,5)\) onto the plane \(x=0\), namely
\[
(0,2,5).
\]
Hence the circle in the \(yz\)-plane has equation
\[
(y-2)^2+(z-5)^2=7.
\]

\medskip
\textbf{Answer:}
\[
\boxed{(x+3)^2+(y-2)^2+(z-5)^2=16}
\]
and the intersection with the \(yz\)-plane is
\[
\boxed{(y-2)^2+(z-5)^2=7}
\]
which is a circle centered at \((0,2,5)\) with radius \(\sqrt{7}\).
\end{solutionbox}

\end{enumerate}

\end{document}
